% ============================================================
%  Chapter: Introduction
% ============================================================
\chapter{Introduction}
\label{chap:introduction}

% ============================================================
\section{Motivation and Context}
\label{sec:motivation}
% ============================================================

The quality of training data is widely recognised as one of the most critical
determinants of the performance and reliability of machine learning
models~\cite{gudivada2017data, whang2023data}.  Real-world datasets are seldom
pristine: they are pervasively affected by a wide range of imperfections
including missing values, mislabelled instances, duplicate records, noisy
measurements, and statistical outliers.  These defects may arise from
heterogeneous sources — sensor malfunction, human annotation error, data
integration across disparate systems, adversarial manipulation, or simply the
inherent noise of the phenomenon being measured.  When such corrupted data are
consumed by learning algorithms, the resulting models may exhibit degraded
generalisation, biased predictions, or — in safety-critical domains such as
cybersecurity and healthcare — potentially harmful decision-making.

Within the paradigm of \emph{Data-Centric AI}~\cite{zha2023datacentric}, the
focus of the machine learning pipeline shifts from model architecture design to
systematic improvement and curation of the training data.  Central to this
paradigm is the ability to \emph{quantify} the impact of specific data quality
issues on downstream model performance.  This capability, however, presupposes
the availability of tools that can introduce controlled, reproducible
perturbations into a dataset while leaving all other aspects of the data intact.
Such tools enable researchers to answer questions of the form: ``How much does
model accuracy degrade when $x$\% of a given feature's values are replaced by
noise?'', or ``Which features are most vulnerable to label corruption?''

\textbf{PuckTrick}\footnote{\url{https://andreamaurino.github.io/pucktrick-ui-docs/}}
is an open-source Python library developed at the \emph{Dat*AI laboratory} of
the University of Milano-Bicocca that addresses precisely this need.  PuckTrick
provides a unified, programmatic interface for the controlled introduction of
five fundamental categories of data corruption into tabular datasets intended for
machine learning applications: \emph{duplicated records}, \emph{missing values},
\emph{noisy values}, \emph{outliers}, and \emph{label inconsistencies}.  Each
corruption type is parametrised by the target column and the desired percentage
of affected rows, enabling fine-grained, reproducible experimental designs in
which individual data quality dimensions can be isolated and studied
independently.

In its original form, PuckTrick operates exclusively on Pandas
DataFrames~\cite{mckinney2010data} — a widely used in-memory data structure for
tabular data in the Python ecosystem.  While Pandas is an excellent choice for
small-to-medium datasets that fit comfortably in a single machine's memory, it
becomes a bottleneck when the dataset scales to millions or tens of millions of
rows.  Modern cybersecurity, IoT, and financial datasets routinely exceed the
single-node memory capacity, and the scientific community increasingly relies on
distributed computing frameworks to process them.

% ============================================================
\section{The Cybersecurity Application Domain}
\label{sec:cybersecurity-domain}
% ============================================================

The empirical component of this thesis is situated within the domain of
\emph{network intrusion detection} — a field that offers a particularly
compelling setting for studying the interplay between data quality and deep
learning performance.

Modern enterprise networks generate vast volumes of traffic, ranging from
routine HTTP/HTTPS exchanges and DNS queries to protocol-level heartbeats and
management flows.  Concealed within this predominantly benign stream, malicious
actors conduct a broad spectrum of attacks — brute-force credential guessing,
Denial-of-Service (DoS) and Distributed Denial-of-Service (DDoS) floods,
web-application exploits (SQL injection, cross-site scripting), botnet
command-and-control communications, and stealthy infiltration campaigns.
A \emph{Network Intrusion Detection System} (NIDS) is tasked with
automatically classifying each observed network flow as either benign or
malicious and, in the latter case, identifying the specific type of attack
being carried out~\cite{aldweesh2020deep}.

This classification problem is inherently challenging for several reasons that
make it an ideal testbed for data quality experimentation:

\begin{enumerate}
    \item \textbf{Massive scale.}  A single day of traffic capture at an
          institutional network can produce millions of flow records, each
          described by dozens of statistical features (packet counts, byte
          volumes, inter-arrival times, TCP flag distributions, etc.).
          Processing and corrupting datasets of this magnitude requires
          distributed computing capabilities — precisely the motivation for
          extending PuckTrick to Apache Spark.
    \item \textbf{Extreme class imbalance.}  In any realistic network
          capture, benign flows vastly outnumber attack flows; moreover, the
          attack categories themselves exhibit highly unequal frequencies
          (volumetric DDoS attacks generate far more flows than low-rate
          infiltrations).  Class imbalance amplifies the sensitivity of
          classifiers to data quality issues: noise injected into minority
          attack classes may have a disproportionate effect on recall and
          F1-score.
    \item \textbf{Absence of payload information.}  For privacy and
          encryption reasons, modern NIDS datasets rely exclusively on
          transport- and network-layer flow statistics rather than
          application-layer payload content.  This constraint places an
          inherent ceiling on discriminative power and makes each feature's
          contribution to classification proportionally more important — and
          its corruption proportionally more damaging.
    \item \textbf{Temporal structure.}  Network attacks are not isolated
          events: they unfold as campaigns with characteristic temporal
          signatures — bursty DDoS floods, periodic port scans, sustained
          brute-force attempts.  Models that exploit this temporal context
          (such as the Temporal Fusion Transformer employed in this work)
          gain a significant advantage, but their temporal reasoning is also
          vulnerable to corruption of the timestamp feature or of the features
          that encode temporal dynamics.
    \item \textbf{Real-world criticality.}  In a production security
          operations centre, a classifier trained on corrupted data may fail
          to detect genuine attacks (false negatives) or overwhelm analysts
          with spurious alerts (false positives).  Understanding how data
          quality degradation maps to these operational failure modes is
          therefore of direct practical relevance.
\end{enumerate}

The specific dataset selected for this study is the \textbf{CSE-CIC-IDS2018}
(\emph{Canadian Institute for Cybersecurity — Intrusion Detection System 2018}),
a publicly available benchmark produced by the University of New
Brunswick~\cite{cicids2018} and distributed on the Kaggle
platform\footnote{\url{https://www.kaggle.com/datasets/solarmainframe/ids-intrusion-csv/data}}.
The dataset was collected over a ten-day capture window in February--March 2018
inside a controlled cyber-range infrastructure: 50 attack machines launched a
curated battery of attacks against a victim network comprising 420 machines and
30 servers, while a profile generator replayed realistic benign traffic drawn
from a usage model encompassing HTTP, HTTPS, FTP, SSH, and email activity.
All bi-directional network flows were extracted by the \textit{CICFlowMeter}
tool~\cite{lashkari2017characterization}, which computes 80 statistical features
— packet-level and flow-level statistics — from raw \texttt{.pcap} captures.
Each flow is labelled as either \texttt{Benign} or one of fourteen mutually
exclusive attack categories:

\begin{itemize}
    \item \emph{Brute-Force} attacks (FTP-BruteForce, SSH-BruteForce);
    \item \emph{Denial-of-Service} (DoS-GoldenEye, DoS-Hulk,
          DoS-SlowHTTPTest, DoS-Slowloris);
    \item \emph{Web attacks} (SQL Injection, XSS, Brute-Force on web login);
    \item \emph{Botnet} traffic (ARES botnet);
    \item \emph{Infiltration} scenarios;
    \item \emph{Distributed Denial-of-Service} (DDoS) attacks with LOIC
          UDP and HOIC.
\end{itemize}

In its unified form the dataset comprises over \textbf{16 million network-flow
instances}, making it one of the largest and most realistic publicly available
benchmarks for intrusion detection research.  A detailed statistical
characterisation of the dataset — including its dimensionality, class
distribution, temporal structure, missing-data profile, feature correlation
structure, distributional properties, and feature fragility classification — is
provided in Chapter~\ref{chap:dataset}.

The choice of this domain is deliberate: by studying the effects of controlled
noise injection on deep learning models that must discriminate between benign
and malicious network traffic, this thesis bridges two active research areas —
\emph{data quality for machine learning} and \emph{AI-driven cybersecurity} —
and provides actionable insights into how resilient (or brittle) modern deep
learning architectures are when the data on which they are trained exhibit
quality degradation of known type and magnitude.

% ============================================================
\section{Objectives of the Thesis}
\label{sec:objectives}
% ============================================================

This thesis pursues two interrelated objectives.

\paragraph{Objective 1: Extending PuckTrick to Apache Spark.}
The primary engineering contribution of this work is the extension of the
PuckTrick library to support \textbf{Apache Spark}~\cite{zaharia2016apache}
(via PySpark) as an alternative processing backend alongside Pandas.  The work
was conducted within the \texttt{Spark\_Conversion} branch of the PuckTrick
repository in collaboration with Daniel Satriano, under the supervision of
Professor Andrea Maurino.

The central design principle guiding the extension was \emph{functional
equivalence}: for every corruption method offered by PuckTrick, the Spark
implementation must accept the same parametrisation and produce statistically
equivalent output distributions, so that experiments conducted on a
single-machine Pandas backend yield results that are directly comparable with
those obtained in a distributed Spark cluster.  This requirement poses
non-trivial engineering challenges, because Pandas and Spark embody
fundamentally different computation models (eager, mutable, row-indexed
DataFrames versus lazy, immutable, partition-based Resilient Distributed
Datasets), and many of PuckTrick's corruption algorithms rely on random
sampling, in-place mutation, and index-based row selection — operations that do
not have direct analogues in the Spark API.

The implementation involved:
\begin{enumerate}
    \item A systematic audit of all existing PuckTrick corruption methods to
          catalogue their Pandas-specific dependencies.
    \item The design of a \emph{dual-backend architecture} in which each
          corruption module detects the type of the input DataFrame (Pandas
          \texttt{pd.DataFrame} or PySpark \texttt{pyspark.sql.DataFrame})
          and dispatches to the appropriate implementation.
    \item The re-implementation of each corruption method using PySpark's
          SQL-based DataFrame API, paying particular attention to
          pseudorandom-number generation with seed control, monotonically
          increasing identifiers, and window functions to simulate
          index-based operations.
    \item An extensive validation campaign comparing the outputs of the Pandas
          and Spark backends across all corruption types, all supported data
          types (continuous, discrete, binary, categorical string,
          categorical integer), and a range of corruption percentages, to
          verify statistical equivalence.
    \item Performance benchmarking on a multi-node Spark cluster to evaluate
          the distributed scalability of the extended library.
\end{enumerate}

\paragraph{Objective 2: Evaluating the impact of controlled noise on deep
learning models for cybersecurity.}
The second objective is an empirical investigation into the effects of
PuckTrick-generated data corruption on the performance of deep learning models
trained for \emph{network intrusion detection}, the cybersecurity application
domain introduced in Section~\ref{sec:cybersecurity-domain}.  This objective
is motivated by a practical question of considerable importance to the security
community: \emph{to what extent can deep learning classifiers, trained on
massive volumes of network-flow data for the purpose of distinguishing benign
from malicious traffic, tolerate degradation in the quality of their training
data?}

The investigation operates at two levels of granularity.  At the
\emph{binary classification} level, the models are asked to discriminate
between benign and malicious flows — a task analogous to anomaly detection,
where any deviation from normal traffic triggers an alert.  At the
\emph{multi-class classification} level, the models must additionally identify
the specific type of attack (e.g., DDoS-Hulk versus SQL Injection versus
Botnet), a substantially harder task that requires the classifier to learn
fine-grained distinctions between attack profiles that often share similar
flow-level statistics.  Studying both levels simultaneously is essential
because data corruption may have asymmetric effects: noise that does not
materially degrade binary detection (benign vs.\ malicious) may nonetheless
collapse the classifier's ability to distinguish among attack subtypes, or
vice versa.

The experimental testbed is the \textbf{CSE-CIC-IDS2018} dataset described in
Section~\ref{sec:cybersecurity-domain}, comprising over 16 million labelled
network-flow instances spanning benign traffic and fourteen distinct attack
categories (detailed in Chapter~\ref{chap:dataset}).  Two state-of-the-art
deep learning architectures were selected for the evaluation:

\begin{itemize}
    \item \textbf{TabNet}~\cite{arik2021tabnet}: a sequential attention-based
          architecture specifically designed for tabular data.  TabNet performs
          instance-wise feature selection through learnable sparse attention
          masks at each decision step, producing natively interpretable feature
          importance scores.  It was configured as a \emph{multi-task
          classifier}, simultaneously predicting both the binary
          (benign/malicious) and the fine-grained multi-class (specific attack
          type) labels.
    \item \textbf{Temporal Fusion Transformer (TFT)}~\cite{lim2021temporal}: a
          transformer-based architecture originally proposed for multi-horizon
          time-series forecasting, here repurposed for sequential flow
          classification.  TFT incorporates \emph{Variable Selection Networks}
          (VSNs) that learn which input features are most informative at each
          time step, and multi-head self-attention over the encoder's temporal
          dimension.  By treating the chronologically ordered network flows as
          a single continuous time series, TFT can capture temporal patterns
          in attack campaigns (e.g., DDoS bursts, periodic port scans) that
          purely tabular models cannot exploit.
\end{itemize}

The experimental protocol is structured as a multi-stage pipeline:

\begin{enumerate}
    \item \textbf{Hyperparameter optimisation on a reduced dataset.}
          A stratified 0.1\% sample of the full dataset is used to conduct
          Bayesian hyperparameter optimisation via Optuna~\cite{akiba2019optuna}
          (Tree-structured Parzen Estimator sampler) for both TabNet and TFT.
          The best hyperparameter configurations are persisted for subsequent
          reuse.
    \item \textbf{Feature importance analysis and selection.}
          Using the baseline (clean) data, the attention-based feature importance
          scores of TabNet and the Variable Selection Network weights of TFT
          are extracted.  The approximately 10--15 most important features for
          each model are identified, and the intersection of these two sets
          yields a compact core of features that are critical for both
          architectures.
    \item \textbf{Feature fragility cross-referencing.}
          The feature importance rankings are cross-referenced with the
          fragility analysis performed during dataset characterisation
          (Chapter~\ref{chap:dataset}), which classifies each feature as high-,
          medium-, or low-fragility based on its skewness and kurtosis.  The
          final set of target features for corruption comprises those that are
          simultaneously \emph{highly important} and \emph{highly fragile} — a
          combination that maximises the expected impact of noise injection on
          model performance.
    \item \textbf{Systematic noise injection and retraining.}
          The dataset is scaled to 5\% of its original size (approximately
          800\,000 instances), retaining only the selected features.  Each of
          the five target features is then corrupted independently using four
          of PuckTrick's five corruption methods (missing, noise, outliers,
          labels), yielding $5 \times 4 = 20$ corrupted dataset variants.  The
          fifth method, \emph{duplicated}, is row-level and therefore applied
          once globally on an arbitrary feature.  Each variant is used to
          retrain both TabNet and TFT, producing a total of
          $(5 \times 4 + 1) \times 2 = 42$ model evaluations.
    \item \textbf{Statistical robustness via repeated trials.}
          The entire corruption-and-retraining loop is repeated 10 times with
          different random seeds to establish confidence intervals around the
          observed performance deltas, ensuring that the results are
          statistically robust and not an artefact of a particular random
          split or noise realisation.
\end{enumerate}

% ============================================================
\section{Technologies and Frameworks}
\label{sec:technologies}
% ============================================================

This section provides a concise overview of the principal technologies and
frameworks employed throughout this thesis, situating each within the broader
landscape of data engineering and deep learning.

% ------------------------------------------------------------
\subsection{Apache Spark and PySpark}
\label{sec:tech-spark}
% ------------------------------------------------------------

Apache Spark~\cite{zaharia2016apache} is a unified analytics engine for
large-scale data processing developed at the AMPLab of UC
Berkeley and subsequently donated to the Apache Software Foundation.  Spark
extends the MapReduce paradigm by introducing an in-memory computation
model based on \emph{Resilient Distributed Datasets} (RDDs) — immutable,
fault-tolerant collections of records partitioned across a cluster's
nodes~\cite{zaharia2012resilient}.  Subsequent releases introduced the
\emph{DataFrame API} (a distributed analogue of the Pandas DataFrame with a
SQL-like query optimiser, Catalyst) and the \emph{Dataset API} (offering
compile-time type safety on the JVM).

PySpark~\cite{spark2024pyspark} is the Python binding to Spark, enabling
Python-native code to orchestrate distributed computation on a Spark cluster.
In the context of this thesis, PySpark provides the execution substrate for the
Spark-based corruption methods implemented in PuckTrick: each corruption
operation is expressed as a sequence of Spark DataFrame transformations —
column selections, filters, joins, window functions, and user-defined functions
(UDFs) — that are compiled by Catalyst into an optimised physical execution plan
and distributed across the cluster's executors.

The Spark cluster used for benchmarking and experimental runs in this work
consists of a driver node and four executor nodes, each equipped with 4 cores
and 13\,GB of memory, for a total distributed capacity of 16 cores and 52\,GB
of RAM.

% ------------------------------------------------------------
\subsection{Pandas}
\label{sec:tech-pandas}
% ------------------------------------------------------------

Pandas~\cite{mckinney2010data} is the \emph{de facto} standard library for
in-memory tabular data manipulation in the Python data science ecosystem.  Its
core abstraction, the \texttt{DataFrame}, is a two-dimensional labelled data
structure with heterogeneous column types, supporting rich indexing, grouping,
merging, and reshaping operations.  Pandas operates eagerly and mutably on a
single machine: all data must reside in RAM, and transformations are executed
immediately upon invocation.

PuckTrick's original implementation relies extensively on Pandas' row-level
indexing (\texttt{.iloc}, \texttt{.loc}), in-place assignment, and
\texttt{numpy}-backed random sampling.  The migration to Spark required
abstracting away these Pandas-specific idioms.

% ------------------------------------------------------------
\subsection{TabNet}
\label{sec:tech-tabnet}
% ------------------------------------------------------------

TabNet~\cite{arik2021tabnet} is a deep learning architecture explicitly designed
for tabular data, proposed by Arik and Pfister at Google Cloud AI.  Its
distinctive mechanism is \emph{sequential attention}: at each of $N_{\text{steps}}$
decision steps, a learnable attention mask selects a sparse subset of input
features to process, and the outputs of successive steps are aggregated to
produce the final prediction.

Key architectural components include:
\begin{itemize}
    \item \textbf{Feature transformer}: a shared-across-steps and
          step-dependent stack of fully connected layers with batch
          normalisation and Gated Linear Units (GLU).
    \item \textbf{Attentive transformer}: a sparsemax~\cite{martins2016softmax}
          or entmax~\cite{peters2019sparse} activated attention mechanism that
          enforces feature sparsity — at each step, only a small number of
          features receive non-zero attention.
    \item \textbf{Sparsity regularisation}: a coefficient $\lambda_{\text{sparse}}$
          penalises the entropy of the attention distribution, encouraging the
          model to rely on a compact set of features.
\end{itemize}

TabNet's attention masks are directly interpretable: by averaging the masks over
all instances and steps, a global feature importance ranking is obtained without
any post-hoc explanation technique.  This property is especially valuable in the
context of noise-injection experiments, as it allows tracking \emph{how} the
model redistributes its attention when specific features are corrupted.

In this work, TabNet is instantiated as a \texttt{TabNetMultiTaskClassifier}
from the \texttt{pytorch-tabnet} library~\cite{dreamquark2020tabnet}, with two
prediction heads sharing a common encoder: one for the binary label
(\texttt{label\_generic}) and one for the multi-class label (\texttt{Label}).
Hyperparameters — including the attention dimension ($N_a = N_d$), the number of
steps ($N_{\text{steps}}$), the relaxation factor ($\gamma$), the sparsity
coefficient ($\lambda_{\text{sparse}}$), the learning rate, batch size, and mask
type — are tuned via Optuna with a Tree-structured Parzen Estimator (TPE)
sampler, optimising the mean Matthews Correlation Coefficient (MCC) across both
tasks.

% ------------------------------------------------------------
\subsection{Temporal Fusion Transformer (TFT)}
\label{sec:tech-tft}
% ------------------------------------------------------------

The Temporal Fusion Transformer~\cite{lim2021temporal} is a
transformer-based architecture developed by Lim \emph{et al.} at Google for
interpretable multi-horizon time-series forecasting.  Its architecture comprises
several specialised sub-networks:

\begin{itemize}
    \item \textbf{Variable Selection Networks (VSNs)}: GRN-based (Gated
          Residual Network) modules that assign learned importance weights to
          each input variable, independently for the encoder (past inputs) and
          the decoder (known future inputs).  VSNs provide per-time-step
          feature importance, enabling the identification of which features are
          informative and \emph{when}.
    \item \textbf{Static covariate encoders}: networks that process
          time-invariant metadata and inject it as context into all
          time-dependent components, enriching the representation with
          entity-level information.
    \item \textbf{Temporal self-attention}: a multi-head attention layer over
          the encoder's temporal dimension, allowing the model to capture
          long-range dependencies across time steps.
    \item \textbf{Gated skip connections and layer normalisation}: used
          throughout the architecture to facilitate gradient flow and prevent
          overfitting.
\end{itemize}

In this thesis, TFT is repurposed as a sequential classification model for
network intrusion detection.  The chronologically ordered network flows are
treated as a single continuous time series, and the model is asked to predict
the labels (binary and multi-class simultaneously via \texttt{MultiLoss}) of the
next flow given an encoder window of the preceding 30 flows.  This formulation
allows TFT to exploit temporal patterns in attack campaigns — for instance,
the bursty nature of DDoS attacks or the sequential probing of port scans — that
purely feature-based models such as TabNet cannot leverage.

Hyperparameters (hidden size, attention head size, dropout rate, hidden
continuous size, learning rate, and batch size) are optimised via Optuna,
minimising the combined cross-entropy validation loss across both prediction
heads.

% ------------------------------------------------------------
\subsection{Optuna}
\label{sec:tech-optuna}
% ------------------------------------------------------------

Optuna~\cite{akiba2019optuna} is an automatic hyperparameter optimisation
framework that implements efficient search strategies including the
\emph{Tree-structured Parzen Estimator} (TPE) and \emph{Hyperband} pruning.
It provides a \emph{define-by-run} API in which the search space is specified
programmatically inside the objective function, enabling dynamic space
definitions that depend on the values of previously sampled hyperparameters.

In this work, Optuna is used to tune both TabNet and TFT.  The TPE sampler is
employed with early-stopping callbacks and configurable trial budgets (10
trials for local development, 20 for cluster-based runs).  Optuna's built-in
visualisation tools (parameter importance plots, optimisation history) are used
to analyse the sensitivity of model performance to individual hyperparameters.

% ------------------------------------------------------------
\subsection{Scikit-learn and PyTorch Ecosystem}
\label{sec:tech-sklearn-pytorch}
% ------------------------------------------------------------

The experimental pipeline relies on several additional frameworks:

\begin{itemize}
    \item \textbf{Scikit-learn}~\cite{pedregosa2011scikit}: used for data
          splitting (\texttt{train\_test\_split} with stratification), feature
          scaling (\texttt{StandardScaler}), label encoding
          (\texttt{LabelEncoder}), and evaluation metrics (accuracy, F1 score,
          Matthews Correlation Coefficient, ROC-AUC, classification report).
    \item \textbf{PyTorch}~\cite{paszke2019pytorch}: the underlying deep
          learning framework for both TabNet (via \texttt{pytorch-tabnet}) and
          TFT (via \texttt{pytorch-forecasting} and
          \texttt{pytorch-lightning}).
    \item \textbf{PyTorch Lightning}~\cite{falcon2019pytorch}: a high-level
          training framework that abstracts the training loop, gradient
          accumulation, checkpointing, and multi-GPU/distributed training,
          used to train the TFT model.
    \item \textbf{PyTorch Forecasting}~\cite{beitner2020pytorchforecasting}:
          a library built on top of PyTorch Lightning that provides pre-built
          time-series architectures including TFT, along with the
          \texttt{TimeSeriesDataSet} abstraction for efficient batched
          sequence generation.
\end{itemize}

% ============================================================
\section{Thesis Structure}
\label{sec:thesis-structure}
% ============================================================

The remainder of this thesis is organised as follows:

\begin{description}
    \item[Chapter~\ref{chap:state-of-the-art}] reviews the state of the art
        in data quality for machine learning and describes the original
        architecture of the PuckTrick library prior to the modifications
        introduced in this work.
    \item[Chapter~\ref{chap:dataset}] presents a comprehensive analysis of the
        CIC-IDS2018 dataset, including its schema, class distribution,
        temporal structure, missing-data profile, feature correlation structure,
        distributional properties, and feature fragility classification.
    \item[Chapter~\ref{chap:spark-extension}] details the design, implementation,
        and validation of the PySpark extension of PuckTrick, including the
        dual-backend architecture, the re-implementation of all five corruption
        modules, and the performance benchmarking results.
    \item[Chapter~\ref{chap:experiments}] describes the experimental protocol
        for evaluating the impact of controlled noise injection on TabNet and
        TFT, presents the results across all corruption types and features, and
        discusses the observed patterns and their implications.
    \item[Chapter~\ref{chap:conclusion}] summarises the contributions of this
        thesis and outlines directions for future work.
\end{description}
